\chapter{Structural Ramsey Theory}

We will now recall some concepts and results from Ramsey theory, for which
we refer the reader to Ne\v{s}et\v{r}il \cite{N95}, Ne\v{s}et\v{r}il-R\"odl
\cite{NR90} and Graham-Rothschild-Spencer \cite{GRS}.

Let $\bfa ,\bfb$ be structures in a signature $L$.  We write
\[
    \bfa\leq\bfb
\]
if $\bfa$ can be embedded into $\bfb$.  If $\bfa\leq\bfb$, we let
\[
    \begin{pmatrix}\bfb\\ \bfa\end{pmatrix}=\{\bfa_0:\bfa_0\text{ is a
        substructure of }\bfb\text{ isomorphic to }\bfa\}.
\]
For $\bfa\leq\bfb\leq\bfc ,\ k=2,3,\dots$, we write, using the Erd\"{o}s-Rado \cite{ER56} arrow-notation,
\[
    \bfc\rightarrow (\bfb )^\bfa_k
\]
if for any coloring
\[
    c:\begin{pmatrix}\bfc\\ \bfa\end{pmatrix}
    \rightarrow\{1,\dots ,k\}
\]
with $k$ colors, there is
\[
    \bfb_0\in
    \begin{pmatrix}\bfc \\ \bfb\end{pmatrix}
\]
which is {\it homogeneous}, in the
\[
    \bfa_0\in\begin{pmatrix}
        \bfb_0\\ \bfa\end{pmatrix},\ c(\bfa_0)=i
\]
sense that for some $1\leq i\leq k$, and all , i.e.,
\[
    \begin{pmatrix}\bfb_0 \\
        \bfa\end{pmatrix}
\]
is monochromatic.

Let now $\k$ be a class of finite structures in a signature $L$.  We
say that $\k$ satisfies the {\it Ramsey property} if $\k$ is hereditary (i.e., satisfies HP) and for any $\bfa
    \leq\bfb$ in $\k ,\ k=2,3,\dots$, there is $\bfc\in\k$ with $\bfb\leq\bfc$
such that
\[
    \bfc\rightarrow (\bfb )^\bfa_k.
\]
Note that by a simple induction, we can restrict this condition to $k=2$.
\begin{mynote}
    Proof of this.
\end{mynote}
Let us now mention some examples of classes with the Ramsey property:

\begin{enumerate}
    \item[(i)] Let $L=\{<\}$ and let $\l\o$ be the class of finite linear orderings. Then $\l\o$ has the Ramsey property, by the classical Ramsey theorem.
          \begin{mynote}
              Proof of this.
          \end{mynote}

    \item [(ii)] Let $L=\{<,E\},\ <,E$ binary relation symbols and let $\o\g$ be the class of all finite ordered graphs $\bfa =\langle A,<^\bfa ,E^\bfa\rangle$ (i.e., $<^\bfa$ is a linear ordering of $A$ and $E^\bfa$ is a symmetric, irreflexive relation).  Then Ne\v{s}et\v{r}il-R\"odl \cite{NR77}, \cite{NR83} showed that $\o\g$ has the Ramsey property.
    \item [(iii)] Let $F$ be a finite field and let $L=\{+\}\cup\{f_\alpha  \}_{\alpha\in F}$ (all function symbols), where $+$ has arity 2 and each $f_\alpha$ is unary. Any vector space over $F$ can be viewed as a structure  in this language with $+$ representing addition and $f_\alpha$ scalar multiplication by $\alpha\in F$.  A substructure of a vector space is clearly a subspace.  Let $\calv_F$ be the class of all finite-dimensional vector spaces over $F$.  Clearly $\bfa\leq\bfb\Leftrightarrow\dim\bfa\leq\dim\bfb$.  It was shown in Graham-Leeb-Rothschild \cite{GLR}, that $\calv_F$ has the Ramsey property.
    \item [(iv)] Let now $L=\{0, 1, - ,\wedge ,\vee\}$ (all function  symbols), where $0, 1$ have arity 0 and $-,\wedge ,\vee$ have arities 1, 2, 2 resp.  Any Boolean algebra is a structure in this language (with $-$ representing Boolean complementation).  Substructures are again subalgebras.  Let $\b\a$ be the class of all finite Boolean algebras.  Then the so-called Dual Ramsey Theorem of Graham-Rothschild \cite{GR} can be  equivalently reformulated by saying that $\b\a$ has the Ramsey property  (see Ne\v{s}et\v{r}il \cite{N95}, 4.13).

\end{enumerate}







Finally, let us point out the following connection between the Ramsey
property and the amalgamation property, discussed in the previous section
(see Ne\v{s}et\v{r}il-R\"odl \cite{NR77}, p. 294, Lemma 1).


\begin{theorem}
    Let $\k$ be a class of finite {\it rigid} (i.e., having no non-trivial automorphisms structures in a signature $L$, which is hereditary.  If $\k$ has the JEP and the Ramsey property, then $\k$ has the AP.
\end{theorem}

\begin{proof}
    To see this, fix $$\bfa ,\bfb ,\bfc\in\k
        \text{ and embeddings } f:\bfa\rightarrow\bfb ,\ g:\bfa\rightarrow \bfc$$  By the
    JEP, find $$\bfe\in\k \text{ in which both } \bfb ,\bfc \text{  can be embedded. }$$
    Then find $$\bfd\in\k \text{ such that } \bfd\rightarrow (\bfe )^\bfa_4$$ and consider
    the coloring $$c:\begin{pmatrix}\bfd\\ \bfa\end{pmatrix}\rightarrow\{x:
        x\subseteq\{\bfb ,\bfc\}\}$$ defined as follows:  Given $$\bfa_0\in\begin{pmatrix}
            \bfd \\ \bfa\end{pmatrix},\bfb\in c(\bfa_0)\Leftrightarrow \text{there is an
            embedding } r:\bfb\rightarrow\bfd \text{ with } r\circ f(\bfa )=\bfa_0$$ and
    similarly for $\bfc$.  Let $$\bfe_0\in\begin{pmatrix}\bfd \\ \bfe
        \end{pmatrix} \text{ be a homogeneous set.}$$   Then $$c(\bfa_0)=\{\bfb ,\bfc\} \text{ for
            all } \bfa_0\in\begin{pmatrix}\bfe_0\\ \bfa\end{pmatrix}$$  For such $\bfa_0$,
    there is $$r:\bfb\rightarrow\bfd  \text{ with } f\circ r(\bfa )=\bfa_0$$ and $s:
        \bfc\rightarrow\bfd$ with $g\circ s(\bfa )=\bfa_0$.  So $r\circ f,\ g\circ s$
    are isomorphisms of $\bfa$ with $\bfa_0$.  Since $\bfa ,\bfa_0$ are rigid,
    it follows that $r\circ f=s\circ g$ so $\bfd ,r,s$ verify the amalgamation
    property for $f:\bfa\rightarrow\bfb ,g:\bfa\rightarrow \bfc$.

    In particular, if $\k$ is a non-empty class of rigid finite structures which is countable, contains structures of arbitrarily
    large cardinality, and satisfies HP and JEP and the Ramsey property, then
    $\k$ is a Fra\"{\i}ss\'{e} class, i.e., the age of a countably infinite
    ultrahomogeneous structure.
\end{proof}

